\chapter{Binary Strings}

Binary strings are the second big family of objects in the enumeration half. Sets of binary
strings can be described compactly by \emph{regular expressions}, and when the expression is
unambiguous it translates directly into a generating series (by length).

\section{Strings and regular expressions}

\begin{definition}[Binary string]
A \emph{binary string} of length $n$ is a sequence $b_1b_2\cdots b_n$ with each $b_i \in\{0,1\}$.
The unique string of length $0$ is the \emph{empty string}, written $\eps$. The
\emph{concatenation} of $a = a_1\cdots a_m$ and $b = b_1\cdots b_n$ is
$ab = a_1\cdots a_m b_1\cdots b_n$.
\end{definition}

For sets of strings $A, B$ define
\begin{itemize}
  \item $AB = \{ab : a \in A, b \in B\}$ (concatenation of sets),
  \item $A^k = AA\cdots A$ ($k$ copies), $A^0 = \{\eps\}$,
  \item $A^* = \bigcup_{k\ge0} A^k$ (the \emph{star} of $A$: all concatenations of zero or more
        strings from $A$).
\end{itemize}
So $\{0,1\}^*$ is the set of all binary strings.

\begin{definition}[Regular expression]
A \emph{regular expression} is defined recursively:
\begin{itemize}
  \item $\eps$, $0$ and $1$ are regular expressions;
  \item if $R$ and $S$ are regular expressions, so are $R \mid S$ (union), $RS$ (concatenation)
        and $R^*$ (star).
\end{itemize}
Each regular expression $R$ \emph{produces} a set of strings $\mathcal{R}$:
$\eps \mapsto \{\eps\}$, $0 \mapsto \{0\}$, $1\mapsto\{1\}$,
$R\mid S \mapsto \mathcal{R}\cup\mathcal{S}$, $RS \mapsto \mathcal{R}\mathcal{S}$,
$R^* \mapsto \mathcal{R}^*$. A set of strings is \emph{regular} if it is produced by some
regular expression.
\end{definition}

\begin{example}
$(0\mid1)^*$ produces all binary strings. $0^*1^*$ produces strings of the form
$00\cdots011\cdots1$. $(0\mid 1)^*11(0\mid1)^*$ produces the strings containing $11$.
\end{example}

\section{Unambiguous expressions}

\begin{definition}[Unambiguous]
A regular expression is \emph{unambiguous} if every string in the set it produces is produced in
exactly one way. Precisely: $R \mid S$ is unambiguous if $\mathcal{R}\cap\mathcal{S} = \emptyset$;
$RS$ is unambiguous if each string of $\mathcal{R}\mathcal{S}$ factors as $rs$
($r \in \mathcal{R}, s \in \mathcal{S}$) in only one way; $R^*$ is unambiguous if each string of
$\mathcal{R}^*$ is a concatenation of strings of $\mathcal{R}$ in only one way; and all
sub-expressions are unambiguous.
\end{definition}

\begin{example}
\begin{itemize}
  \item $(0\mid 00)^*$ is ambiguous: $00$ is produced as $0\cdot0$ and as $00$.
  \item $(0\mid1)^*\mid 0^*$ is ambiguous: $00$ is in both parts of the union.
  \item $(0\mid 01)^*$ is unambiguous: reading a string from the right, a $1$ must be the end of
        a $01$ and a $0$ not followed by $1$ must be a lone $0$, so the factorization is forced.
  \item $(0\mid1)^*11(0\mid1)^*$ is ambiguous: $111$ factors as $\eps\cdot11\cdot1$ and as
        $1\cdot11\cdot\eps$.
\end{itemize}
\end{example}

\begin{theorem}[Sum, product and star lemmas for strings]
Let $R, S$ be regular expressions, and weight strings by length. If $R\mid S$, $RS$ and $R^*$
are unambiguous then
\[
\Phi_{R\mid S} = \Phi_R + \Phi_S,\qquad \Phi_{RS} = \Phi_R\Phi_S,\qquad
\Phi_{R^*} = \frac{1}{1 - \Phi_R},
\]
and $\Phi_\eps = 1$, $\Phi_0 = \Phi_1 = x$.
\end{theorem}
\begin{proof}
Unambiguity says exactly that the maps $\mathcal{R}\times\mathcal{S}\to\mathcal{R}\mathcal{S}$,
$(r,s)\mapsto rs$ and $\bigcup_k \mathcal{R}^k \to \mathcal{R}^*$ are bijections, and that the
union is disjoint. Length is additive under concatenation, so these are the sum, product and
string lemmas of the previous chapter. (For the star, $\eps \notin \mathcal{R}$ is forced by
unambiguity, since otherwise $\eps = \eps\eps$ has two factorizations.)
\end{proof}

\begin{important}[Pitfall]
Applying these rules to an ambiguous expression gives a series that \emph{overcounts}. E.g.\
$(0\mid 00)^*$ would give $\frac{1}{1-x-x^2}$, but the set it produces is $0^*$ with series
$\frac{1}{1-x}$. Always justify unambiguity before translating.
\end{important}

\begin{example}
$\Phi_{(0\mid1)^*} = \frac{1}{1-2x}$, so there are $2^n$ strings of length $n$.
$(0\mid11)^*$ (strings built from tokens $0$ and $11$) has series $\frac{1}{1-x-x^2}$, and
$(0 \mid 01)^*$ also has series $\frac{1}{1-x-x^2}$.
\end{example}

\section{Block decompositions}

\begin{definition}[Block]
A \emph{block} of a binary string is a maximal non-empty substring of equal bits. For example
$0011101$ has blocks $00$, $111$, $0$, $1$.
\end{definition}

Every binary string is uniquely of the form ``some $0$s, then alternating blocks, then some $1$s'':

\begin{theorem}[Block decompositions]
The expressions
\[
0^*(1\,1^*\,0\,0^*)^*\,1^* \qquad\text{and}\qquad 1^*(0\,0^*\,1\,1^*)^*\,0^*
\]
are unambiguous and both produce all binary strings.
\end{theorem}
\begin{proof}[Proof sketch]
For the first: the leading $0^*$ is the initial block of $0$s (possibly empty); each $11^*00^*$
is a block of $1$s followed by a block of $0$s; the trailing $1^*$ is a final block of $1$s
(possibly empty). Given a string, its blocks determine each piece uniquely.
\end{proof}

Check: $\Phi = \frac{1}{1-x}\cdot\frac{1}{1 - \left(\frac{x}{1-x}\right)^2}\cdot\frac{1}{1-x}
= \frac{1}{(1-x)^2 - x^2} = \frac{1}{1-2x}$, as it must be.

The point of block decompositions is that restrictions on block lengths are easy to impose:
replace $1\,1^*$ by an expression for the allowed block lengths.

\begin{example}[Every block of 1s has even length]
Replace $11^*$ by $11(11)^*$ (blocks of even positive length) and $1^*$ at the end by
$\eps \mid 11(11)^*$:
\[
0^*\big(11(11)^*\,0\,0^*\big)^*\big(\eps\mid 11(11)^*\big).
\]
With $A = \frac{1}{1-x}$, $B = \frac{x^2}{1-x^2}$ and $C = \frac{x}{1-x}$ (for $00^*$), and noting
$1 + B = \frac{1}{1-x^2}$,
\[
\Phi = \frac{1}{1-x}\cdot\frac{1}{1 - \frac{x^2}{1-x^2}\cdot\frac{x}{1-x}}\cdot\frac{1}{1-x^2}
= \frac{1}{(1-x)(1-x^2) - x^3} = \frac{1}{1-x-x^2}.
\]
Indeed $(1-x)(1-x^2) = 1 - x - x^2 + x^3$. This agrees with the simpler description $(0\mid11)^*$.
Check $n = 3$: $000$, $110$, $011$; and $\coef{3}\frac{1}{1-x-x^2} = 3$.
\end{example}

\begin{example}[No block of length more than 2]
Allowed blocks of 1s: $1 \mid 11$, series $x + x^2$; same for 0s. Using
$(\eps\mid 0\mid00)\big((1\mid11)(0\mid00)\big)^*(\eps\mid1\mid11)$:
\[
\Phi = \frac{(1+x+x^2)^2}{1-(x+x^2)^2}.
\]
Its coefficients are $1,2,4,6,10,\dots$; for $n=3$ the forbidden strings are $000$ and $111$,
leaving $6$. (Expanding: $(1 + x+x^2)^2 = 1 + 2x + 3x^2 + 2x^3 + x^4$ and $(x+x^2)^2 = x^2 + 2x^3
+ x^4$, so $a_n = a_{n-2} + 2a_{n-3} + a_{n-4} + p_n$, giving $a_0=1$, $a_1 = 2$, $a_2 = 1 + 3 = 4$,
$a_3 = 2 + 2 + 2 = 6$, $a_4 = 4 + 4 + 1 + 1 = 10$.)
\end{example}

\section{Prefix decompositions}

A \emph{prefix decomposition} cuts a string after each occurrence of some marker, describing
the string as a sequence of units that each end in the marker, followed by a tail.

\begin{example}
$\{0,1\}^* = (1^*0)^*1^*$ unambiguously: cut after every $0$. Each unit is a (possibly empty)
run of $1$s followed by a $0$, and the tail is the final run of $1$s. Check:
$\frac{1}{1 - \frac{x}{1-x}}\cdot\frac{1}{1-x} = \frac{1}{1-2x}$.
\end{example}

\begin{example}[No $1$ is followed by another $1$]
Strings avoiding $11$: cut after each $0$; each unit is $0$ or $10$, and the tail is $\eps$ or
$1$. So the expression is $(0\mid10)^*(\eps\mid1)$, with series $\frac{1+x}{1-x-x^2}$:
coefficients $1, 2, 3, 5, 8, \dots$.
\end{example}

For an arbitrary forbidden substring $\kappa$ there is a general method. Let $A$ be the set of
strings that avoid $\kappa$ and $B$ the set of strings that contain $\kappa$ exactly once, as a
suffix. Then
\begin{enumerate}
  \item $\{\eps\} \cup A\{0,1\} = A \cup B$ (disjointly): appending one bit to a string of $A$
        either still avoids $\kappa$, or creates exactly one copy at the end. Hence
        $1 + 2xA(x) = A(x) + B(x)$.
  \item $A\{\kappa\} = B \cup \bigcup B\{\text{leftover}\}$, one term for each proper suffix of
        $\kappa$ that is also a prefix of $\kappa$ (an \emph{overlap}): the first occurrence of
        $\kappa$ in $a\kappa$ may end early, inside the appended $\kappa$.
\end{enumerate}

\begin{example}[Avoiding $101$]
The only overlap is the suffix $1$ (equal to the prefix $1$), leaving $01$. So
$A\{101\} = B \cup B\{01\}$: $x^3A = B(1 + x^2)$. Then
\[
1 + 2xA = A + \frac{x^3A}{1+x^2} \ \Longrightarrow\ A(x) = \frac{1+x^2}{1 - 2x + x^2 - x^3}.
\]
The coefficients satisfy $a_n = 2a_{n-1} - a_{n-2} + a_{n-3}$ for $n \ge 3$ with $a_0 = 1$,
$a_1 = 2$, $a_2 = 4$, giving $a_3 = 7$ and $a_4 = 12$. Check $n=4$: the strings containing
$101$ are $1010, 1011, 0101, 1101$ (no length-4 string has two occurrences), so
$16 - 4 = 12$.
\end{example}

\section{Recursive decompositions}

A set of strings can also be described by a recursive equation, which translates into an
equation for its generating series.

\begin{example}
$S = \{0,1\}^*$ satisfies $S = \{\eps\} \cup \{0\}S \cup \{1\}S$ (disjointly, by the first bit),
so $\Phi_S = 1 + 2x\Phi_S$ and $\Phi_S = \frac{1}{1-2x}$.
\end{example}

\begin{example}[Strings with no $11$, recursively]
Let $S$ be the strings avoiding $11$. A string in $S$ is $\eps$, or $1$, or $0$ followed by a
string in $S$, or $10$ followed by a string in $S$:
\[
S = \{\eps\}\cup\{1\}\cup\{0\}S\cup\{10\}S,\qquad
\Phi_S = 1 + x + x\Phi_S + x^2\Phi_S,\qquad \Phi_S = \frac{1+x}{1 - x - x^2},
\]
matching the prefix decomposition.
\end{example}

\begin{example}[Equal numbers of 0s and 1s, by first return]
Recursive decompositions also handle non-regular sets. Let $D$ be strings with equal numbers of
$0$s and $1$s that never have more $1$s than $0$s in a prefix (``Dyck words''). Then
$D = \{\eps\}\cup\{0\}D\{1\}D$, giving $\Phi_D(x) = 1 + x^2\Phi_D(x)^2$. Solving the quadratic
gives $\Phi_D = \frac{1 - \sqrt{1-4x^2}}{2x^2}$, whose coefficients are the Catalan numbers
$\frac{1}{m+1}\binom{2m}{m}$ at $x^{2m}$.
\end{example}

\section{Summary}
\begin{itemize}
  \item Regular expressions: $\eps,0,1$, union, concatenation, star.
  \item Unambiguous expressions translate: $\mid\ \to +$, concatenation $\to \times$,
        $R^* \to 1/(1-\Phi_R)$. Ambiguous ones overcount.
  \item Block decomposition $0^*(11^*00^*)^*1^*$ for block-length restrictions.
  \item Prefix decomposition $(1^*0)^*1^*$; general substring avoidance via the $A,B$ equations.
  \item Recursive decompositions give functional equations for $\Phi$.
\end{itemize}
